\bookStart{Introduction to Norse Mythic Poetry}

The present section contains almost all extant narrative poetry concerning the pre-Christian Germanic gods, most of which is found in the first half of \Regius, and all of which is in Old Norse.
Of all genres of Germanic poetry, the mythological is in fact the one of which we have the earliest mention.  Tacitus, in \emph{Germania}, ch. 2, talking about the ancient Germani, says that: \textlatin{Celebrant carminibus antiquis, quod unum apud illos memoriae et annalium genus est, Tuistonem deum terra editum.} ‘They celebrate in ancient songs, which is the only kind of memory and history among them, the earth-born god Tuisto.’
In spite of this fact it is not due to any decision on the part of the editor that this type of poetry is exclusively in Old Norse, for outside of Iceland there survive no examples of this undoubtedly ancient genre.
This must obviously not be taken to say that no such poetry existed, something disproven by the Tacitean quote above.  In fact, as we shall see, many of the myths and gods dealt with in this genre are attested throughout the Germanic speech-area for a period of many centuries; in addition they show deep parallels to religious traditions as distant as Mesopotamia, India, and Egypt.  And how should they have been told and transmitted through unreckoned illiterate ages if not in poetic form?

\section{The poems}

The present section includes fourteen poems and eight smaller fragments.

\subsection{Short summaries}

\emph{Vǫlu spǫ́} (‘Spae of the Wallow’, \textlink{Voluspa}) is the prophecy of a wallow, a Germanic seeress, brought back from the dead by the chief god Weden.  The wallow narrates the events of the world’s creation and golden age, its degeneration and ultimate destruction, and at last its rebirth.  The poem is mystical and contains numerous allusions to important myths; it serves as the major source for Snorre’s \Gylfaginning.

\emph{Hǫ́va mǫ́l} (‘Speeches of the High One’, \textlink{Havamal}) is the god Weden’s wisdom, consisting of practical life advice about friends, wealth, wisdom and women, two romantic escapades, and lastly stanzas about runes and galders.  It is by far the longest poem in \Regius.

\emph{Vaf·þrúðnis mǫ́l} (‘Speeches of Webthrithner’, \textlink{Vafthrudnismal}) describes Weden’s visit to the ettin Webthrithner and the ensuing wisdom-contest between them as they wager their heads.  The two ask each other numerous mythological and cosmological questions and in the mouth of Weden these soon turn to the subject of the end times; expectedly the god eventually comes out on top with an unknowable question.

\emph{Grímnis mǫ́l} (‘Speeches of Grimner’, \textlink{Grimnismal}) follows Weden as he sets out to test the hospitality of his former protégé King Garfrith as part of a wager with his wife Frie.  The king is less than friendly and unwittingly tortures the god between two fires; Weden is however relieved by the king’s son, and in exchange gives him a great deal of gnomic wisdom about the dwellings of the gods, sacred locations and objects, the creation of the world, and finally his own nature.  At the end the king dies by his own sword and his son is made king after him.

\emph{Baldrs draumar} (‘Dreams of Balder’, \textlink{Baldrsdraumar}) is a short narrative poem which sees Weden set out to find answers about why his son Balder is having ominous dreams; he brings a wallow back to life and asks her a few questions.

\emph{Skírnis mǫ́l} (‘Speeches of Shirner’, \textlink{Skirnismal}) sees the lovesick god Free send his servant Shirner in order to win him the hand of the beautiful ettin-woman Gird; he does so by means of a curse.

\emph{Hár·barðs ljóð} (‘Leeds of Hoarbeard’, \textlink{Harbardsljod}) is a long argument between the god Thunder and his father Weden, disguised as the ferryman Hoarbeard.  The two stand on opposite sides of a sound and Thunder begs to be let over; he is, of course, refused.

\emph{Hymis kviða} (‘Lay of Hymer’, \textlink{Hymiskvida}) is a humorous poem that tells how the god Thunder and his companion Tew set out to find a huge cauldron for the gods to drink out of.  The two find themselves on a strange adventure as they go fishing with the cauldron’s owner Hymer and Thunder nearly pulls up the Middenyardswyrm in what was undoubtedly one of the most popular myths of the Wiking Age.

\emph{Loka sęnna} (‘Flyting of Lock’, \textlink{Lokasenna}) sees the evil god Lock berate the other gods at a feast, giving personal and often quite graphic insults to each one and alluding to a number of less favourable myths; at last he is banished, caught and placed to be tortured until the end of days like a Northern \textgreek{Προμηθεύς}.

\emph{Þryms kviða} (‘Lay of Thrim’, \textlink{Thrymskvida}) is another humorous poem about Thunder where the god’s hammer is stolen by the ettin Thrim; in order to retrieve it he has to dress up as the goddess Frow and pretend to marry the ettin, yet the poem ends with a great killing as soon as the hammer is brought out.

\emph{Vǫlundar kviða} (‘Lay of Wayland’, \textlink{Volundarkvida}) is properly a heroic poem, but is placed in this section as it is in \Regius.  The poem tells the story of the masterful supernatural smith Wayland, who is kidnapped and hamstrung by the tyrant Nithad and forced to make jewels for him and his family.  Wayland takes his gruesome revenge by killing the tyrant’s sons, impregnating his daughter, Beadhild, and flying away in a flight-suit of his own design.

\emph{Al·víss mǫ́l} (‘Speeches of Allwise’, \textlink{Alvissmal}) is a short poem where a dwarf asks for the hand of Thunder’s daughter.  Thunder tricks him that he will win her if he can only tell him the various poetic names for things (beer, sun, clouds, and the like); the dwarf does so, but in his foolish lust continues for so long that eventually the sun rises and he turns to stone.

\emph{Rígs þula} (‘Thule of Righ’, \textlink{Rigsthula}) narrates the fathering of the three social castes---the thralls, churls, and earls---and the eventual establishment of kingship and elite runic knowledge at the hands of the mysterious god Righ.  The poem is unfortunately incomplete.

\emph{Hyndlu ljóð} (‘Leeds of Hindle’, \textlink{Hyndluljod}) describes how the goddess Frow asks the ettin-woman Hindle to reveal the genealogy of Frow’s lover Oughter so that he will be able to claim his legitimate inheritance.  The poem takes a mythological turn as Hindle briefly recounts the coming destruction of the world.

Following these poems, I offer 8 Eddic fragments (containing a total of 11 stanzas or half-stanzas) cited in \Gylfaginning\ and \Skaldskaparmal\ under the section Fragments from Snorre’s Edda (abbrev.~\textlink{EddicFragments}).  The fragments are edited along with their surrounding prose contexts.  Many of them have likely belonged to now-lost Eddic poems, and their inclusion completes the mythological Eddic corpus.

\subsection{Types of poetry}

The above-listed poems may be categoried on a scale between narrative poetry on the one end and gnomic or wisdom poetry on the other.  Although almost every poem has some type of narrative progression, a distinction can still be made based on the nature of the bulk of the content and whether the narrative exists merely as a frame for the gnomic material or whether it fills an important purpose in its own right.  The exact categorisation is subjective, but for our purposes it may be said that on the far end of the narrative side stand \textlink{Thrymskvida} and \textlink{Hymiskvida}, and on the far end of the gnomic side \textlink{Havamal} and \textlink{Alvissmal}.

\subsection{Exclusions}

A few Eddic-style poems which one the reader may be familiar with from other editions have been excluded from the present edition due to their questionable and likely post-mediæval provenance.  It must be stressed that none of the excluded texts should be relied upon for knowledge of the authentic pre-Christian culture and religion.

The main exclusion when compared with most other modern editions is the antiquarian \emph{Svip·dags mǫ́l} ‘Speeches of Swebday’, which properly consists of the two poems \emph{Fjǫl·svinns mǫ́l} and \emph{Gróu galdr} (\Grougaldr).
Although found only in post-reformation Icelandic paper mss., both parts probably date to the antiquarian period of the 13th century.  Their late origin is proven by the fact that they contain allusions to Christianity and post-conversion language (e.g., \emph{Fjǫl·svinns mǫ́l}: 46/3 \emph{jartegn} ‘evidence’, \Grougaldr\ 13/3: \emph{kristin dauð kona} ‘a dead Christian woman’).
The mythological references they do contain appear to be largely fabricated or derivative of earlier texts, and are likely to introduce much confusion about the mythology on the part of the reader.  That they have recently been translated by \textcite{PettitEdda} further justifies their exclusion.

Apart from \emph{Svip·dags mǫ́l}, two other exclusions found in some older, though not in most modern eds., also merit discussion since the reader may be surprised at their absence.  The \emph{Sólar-ljóð} ‘Leeds of the Sun’ is clearly a late, Christian poem, and is edited and translated by Carolyne Larrington and Peter Robinson in \Skp\ 7.  The \emph{Hrafna-galdr Óðins} ‘Raven-Galder of Weden’ is undoubtedly a post-reformation pastiche as seen by its late meter and use of the Low German loanword \emph{mál-tíð} ‘meal time’.

\section{Table of manuscripts}

The following table shows the manuscript attestation for the included poems and their relative position in each manuscript.  A \emph{q} in \Gylfaginning/\Skaldskaparmal\ represents one or more stanzas attested in quotation.  For more precise details see the Introductions to the individual poems.

{\small \begin{longtable}{|l c c c c c c|}
	\hline
	Signum & \Regius & \AM & \Wormianus & \Hauksbok & \FlatMS & \Gylfaginning/\Skaldskaparmal \\ [0.5ex]
	\hline\hline\endhead
	\hline\endfoot
  \textlink{Voluspa}        & 1  & − & − & + & − & \emph{q} \\
  \textlink{Havamal}        & 2  & − & − & − & − & \emph{q} \\
  \textlink{Vafthrudnismal} & 3  & 4 & − & − & − & \emph{q} \\
  \textlink{Grimnismal}     & 4  & 5 & − & − & − & \emph{q} \\
  \textlink{Skirnismal}     & 5  & 3 & − & − & − & \emph{q} \\
  \textlink{Harbardsljod}   & 6  & 1 & − & − & − & − \\
  \textlink{Hymiskvida}     & 7  & 6 & − & − & − & − \\
  \textlink{Lokasenna}      & 8  & − & − & − & − & \emph{q} \\
  \textlink{Thrymskvida}    & 9  & − & − & − & − & − \\
  \textlink{Volundarkvida}  & 10 & 7 & − & − & − & − \\
  \textlink{Alvissmal}      & 11 & − & − & − & − & \emph{q} \\
  \textlink{Baldrsdraumar}  & −  & 2 & − & − & − & − \\
  \textlink{Rigsthula}      & −  & − & + & − & − & − \\
  \textlink{Hyndluljod}     & −  & − & − & − & + & \emph{q} \\ [1ex]
  \hline
\end{longtable}}
